% !PRINTED 415
\weberpage{415}
Setzt man noch
\weq{6}{a+\partial = y,}
so folgt:
\weq{7}{\begin{aligned}
a&=\frac{y+B\,x}{2},\\
b&=A\,x,\\
c&=-C\,x,\\
\partial&=\frac{y-B\,x}{2},
\end{aligned}}
woraus sich, wenn
\[
a\,\partial - b\,c = n
\]
gesetzt wird, ergibt:
\weq{8}{4\,n = y^2 - D\,x^2.}
Die Zahl $n$ wird also in die beiden komplexen ganzzahligen Faktoren
\[
n = \frac{y+x\sqrt{D}}{2}\cdot\frac{y-x\sqrt{D}}{2}
\]
zerlegt. Für $\omega$ und $\mu$ findet sich noch:
\[
\omega = \frac{-B+\sqrt{D}}{2\,A} = \frac{2\,C}{-B-\sqrt{D}},
\]
\[
\mu = a+b\,\omega = \frac{y+\sqrt{D}\,x}{2}.
\]

Die hier eingeführten Zahlen $x$, $y$ sind nur an die eine Bedingung
geknüpft, daß $y+B\,x$ und folglich auch $y-B\,x$ gerade Zahlen seien,
damit $a$, $\partial$ nach (7) ganze Zahlen werden. Ist also $B$
gerade, so muß $y$ gerade sein, während $x$ beliebig ist; ist $B$
ungerade, so sind $x$ und $y$ beide gerade oder beide ungerade
anzunehmen. Nach (4) kann man diese Unterscheidung auch so ausdrücken:
\[
\begin{aligned}
\text{1) Ist }D&\equiv0\ (\mathrm{mod}\ 4), &\text{so ist }y&\equiv0\ (\mathrm{mod}\ 2),\\
\text{2) Ist }D&\equiv1\ (\mathrm{mod}\ 4), &\text{so ist }y&\equiv x\ (\mathrm{mod}\ 2).
\end{aligned}
\]

Sollen bei gegebenen $A$, $B$, $C$ die Zahlen $a$, $b$, $c$, $\partial$
ohne gemeinsamen Teiler sein, so kommen noch andere Bedingungen hinzu:

Aus (5), (6) folgt, daß jeder gemeinsame Teiler von $a$, $b$, $c$,
$\partial$ auch Teiler von $x$ und $y$ ist; umgekehrt ist jeder
gemeinsame Teiler von $x$ und $y$ auch Teiler von $b$, $c$, $2\,a$,
$2\,\partial$. Sollen also $a$, $b$, $c$, $\partial$ ohne gemeinsamen
Teiler sein, so können $x$ und $y$
%CONT>>
